\subsection*{File 05G: Median/MAD Defence against SGA Poisoning in FL and SFL}

File 05G introduces the first defence mechanism against the poisoning attacks introduced in File 05F. The objective of this experiment was to investigate whether a lightweight robust aggregation defence could stabilise both Federated Learning (FL) and Split Federated Learning (SFL) under the same Scale Gradient Attack (SGA) conditions.

The experiment continues using:
\begin{itemize}
    \item the MNIST dataset,
    \item the extreme one-label-per-client non-IID configuration,
    \item the three-layer MLP architecture,
    \item and the same SGA poisoning attack from File 05F.
\end{itemize}

The model architecture remains:

\[
784 \rightarrow 256 \rightarrow 128 \rightarrow 10
\]

with the SFL split defined as:

\[
\text{Client side: } 784 \rightarrow 256
\]

\[
\text{Server side: } 256 \rightarrow 128 \rightarrow 10
\]

The major change introduced in this notebook is the addition of a Median/MAD-based robust defence mechanism.

\subsubsection*{Why Median was Selected Instead of Mean}

Traditional Federated Averaging uses the arithmetic mean:

\[
\theta^{t+1}
=
\frac{1}{K}
\sum_{k=1}^{K}
\theta_k^{t+1}
\]

However, the arithmetic mean is highly sensitive to outliers. A single malicious client producing extremely large poisoned updates can dominate the aggregation process.

The median operator was selected because it is substantially more robust against outliers. Instead of averaging updates directly, the defence computes the coordinate-wise median across client updates:

\[
m_j
=
\text{median}
\left(
u_{1,j},
u_{2,j},
\dots,
u_{K,j}
\right)
\]

where:
\begin{itemize}
    \item $u_{k,j}$ is parameter coordinate $j$ from client $k$.
\end{itemize}

The defence then measures how far each client deviates from the median update.

This approach was selected because:
\begin{itemize}
    \item the project uses only 10 clients,
    \item and many advanced FL defences such as Krum, Multi-Krum, or Bulyan perform poorly or become unstable under small client counts.
\end{itemize}

Many established FL defence algorithms assume:
\begin{itemize}
    \item large numbers of clients,
    \item or large proportions of honest participants.
\end{itemize}

Under extreme non-IID distributions with only 10 clients:
\begin{itemize}
    \item naturally different client updates already appear highly anomalous,
    \item making neighbour-based FL defences unreliable.
\end{itemize}

Median/MAD was therefore selected because it:
\begin{itemize}
    \item is lightweight,
    \item robust to extreme values,
    \item computationally simple,
    \item and does not require pairwise client-distance calculations.
\end{itemize}

\subsubsection*{MAD-Based Anomaly Detection}

The defence uses the Median Absolute Deviation (MAD) statistic to identify suspicious updates.

First, the coordinate-wise median update is computed:

\[
m
=
\text{median}(u_1, u_2, \dots, u_K)
\]

Then the deviation of each client from the median is calculated:

\[
d_k
=
\|u_k - m\|
\]

The MAD statistic is then:

\[
MAD
=
\text{median}
\left(
|d_k - \text{median}(d)|
\right)
\]

Each client receives a robust anomaly score:

\[
z_k
=
\frac{|d_k - \text{median}(d)|}{MAD + \epsilon}
\]

A client is flagged if:

\[
z_k > \tau
\]

where:
\begin{itemize}
    \item $\tau$ is the MAD threshold coefficient.
\end{itemize}

In this notebook:

\[
\tau = 6
\]

The value was selected experimentally through repeated trial-and-error testing. Smaller thresholds caused:
\begin{itemize}
    \item honest clients to be removed excessively under extreme non-IID conditions.
\end{itemize}

Larger thresholds:
\begin{itemize}
    \item failed to suppress the poisoning sufficiently.
\end{itemize}

The selected value therefore represented a practical balance between:
\begin{itemize}
    \item robustness,
    \item and preservation of legitimate client diversity.
\end{itemize}

\subsubsection*{How Defence Works in FL}

In FL, the malicious client poisons all three layers of the model.

Therefore, the FL defence evaluates:
\begin{itemize}
    \item the full model update vector,
    \item across all layers simultaneously.
\end{itemize}

The malicious client update is completely excluded from aggregation if flagged:

\[
\mathcal{K}_{keep}
=
\{
k \mid z_k \leq \tau
\}
\]

The robust FL aggregation becomes:

\[
\theta^{t+1}
=
\frac{1}{|\mathcal{K}_{keep}|}
\sum_{k \in \mathcal{K}_{keep}}
\theta_k^{t+1}
\]

Thus, flagged FL clients contribute nothing to the next global model.

\subsubsection*{How Defence Works in SFL}

The SFL defence behaves differently because the attack only affects:
\begin{itemize}
    \item the client-side layer.
\end{itemize}

The server-side layers are not directly poisoned by the malicious client.

The SFL defence therefore evaluates:
\begin{itemize}
    \item server-observable gradients,
    \item client-side parameter updates,
    \item and server-side optimisation behaviour.
\end{itemize}

Importantly:
\begin{itemize}
    \item the SFL server never directly accesses client raw data,
    \item preserving privacy.
\end{itemize}

Instead, the server analyses:
\begin{itemize}
    \item activations,
    \item gradients,
    \item and update statistics.
\end{itemize}

This allows anomaly detection without reconstructing private client datasets.

The SFL client-side FL aggregation still occurs separately:

\[
\theta_c^{t+1}
=
\sum_{k=1}^{K}
\frac{n_k}{n}
\theta_{c,k}^{t+1}
\]

while server-side parameters are aggregated independently:

\[
\theta_s^{t+1}
=
\sum_{k=1}^{K}
\frac{n_k}{n}
\theta_{s,k}^{t+1}
\]

The reconstructed SFL model is then formed from both aggregated components.

\subsubsection*{Flagged Client Diagnostics}

The notebook records:
\[
\texttt{FL flagged clients}
\]

and

\[
\texttt{SFL flagged clients}
\]

These diagnostics indicate which clients were classified as anomalous during each communication round.

This became extremely useful during debugging and hyperparameter tuning because:
\begin{itemize}
    \item false positives could be identified,
    \item aggressive thresholds could be detected,
    \item and defence instability could be analysed across rounds.
\end{itemize}

The logs also allowed identification of:
\begin{itemize}
    \item rounds where the malicious client escaped detection,
    \item or rounds where too many honest clients were removed.
\end{itemize}

\subsubsection*{Maximum Absolute Weight Monitoring}

The experiment also records:

\[
\texttt{SFL max abs weight}
\]

and corresponding FL maximum absolute weights.

This metric monitors:

\[
\max_i |\theta_i|
\]

across all model parameters.

This diagnostic was important because poisoning attacks often create:
\begin{itemize}
    \item exploding parameter values,
    \item numerical instability,
    \item NaNs,
    \item or optimisation divergence.
\end{itemize}

Monitoring maximum absolute weights helped identify:
\begin{itemize}
    \item early instability,
    \item exploding gradients,
    \item and ineffective defence thresholds.
\end{itemize}

\subsubsection*{FL Minus SFL Accuracy}

The metric:

\[
A_{diff}
=
A_{FL}
-
A_{SFL}
\]

was retained to compare robustness between systems.

Negative values indicate:
\begin{itemize}
    \item SFL outperforming FL.
\end{itemize}

Positive values indicate:
\begin{itemize}
    \item FL outperforming SFL.
\end{itemize}

This became useful for analysing:
\begin{itemize}
    \item whether the defence stabilised one system more effectively than the other.
\end{itemize}

\subsubsection*{Role of the Global L2 Distance}

The global L2 distance was retained from earlier experiments:

\[
L_2
=
\left\|
\theta_{FL}
-
\theta_{SFL}
\right\|_2
\]

Although FL and SFL are no longer expected to remain equivalent under attack conditions, the L2 metric still provides useful information:
\begin{itemize}
    \item small L2 growth suggests similar optimisation trajectories,
    \item while large L2 divergence indicates fundamentally different optimisation behaviour.
\end{itemize}

Under the defence setting:
\begin{itemize}
    \item the L2 distance no longer represents equivalence validation,
    \item but instead measures behavioural divergence between defended FL and defended SFL systems.
\end{itemize}

\subsubsection*{Results and Interpretation}

Figure~\ref{fig:05g_accuracy} shows a major improvement compared to the undefended attack experiment in File 05F.

The FL model no longer collapses catastrophically. Instead:
\begin{itemize}
    \item both FL and SFL continue stable learning,
    \item and both systems reach accuracies above 0.60.
\end{itemize}

Interestingly:
\begin{itemize}
    \item SFL often slightly outperforms FL during many communication rounds.
\end{itemize}

This suggests that:
\begin{itemize}
    \item the isolated attack surface in SFL,
    \item combined with robust aggregation,
\end{itemize}

creates a naturally stabilised optimisation process.

\begin{figure}[H]
    \centering
    \includegraphics[width=0.85\linewidth]{figures/05G_accuracy.png}
    \caption{FL vs reconstructed SFL global test accuracy using the Median/MAD defence under SGA poisoning attacks. Both systems remain stable and avoid catastrophic collapse.}
    \label{fig:05g_accuracy}
\end{figure}

Figure~\ref{fig:05g_loss} shows that:
\begin{itemize}
    \item the FL loss remains unstable and oscillatory,
    \item while the SFL loss remains comparatively smoother.
\end{itemize}

This behaviour indicates that:
\begin{itemize}
    \item although the defence successfully prevents collapse,
    \item FL still experiences stronger optimisation disruption because the attacker poisons the full model.
\end{itemize}

SFL remains more stable because:
\begin{itemize}
    \item only the client-side layer is attacked,
    \item while the server-side optimisation remains comparatively protected.
\end{itemize}

\begin{figure}[H]
    \centering
    \includegraphics[width=0.85\linewidth]{figures/05G_loss.png}
    \caption{FL vs reconstructed SFL global test loss using the Median/MAD defence. FL remains noisier and less stable, while SFL maintains smoother optimisation behaviour.}
    \label{fig:05g_loss}
\end{figure}
\begin{figure}[H]
    \centering
    \includegraphics[width=0.4\linewidth]{figures/median_mad_defence.png}
    \caption{Median/MAD poisoning defence workflow.}
    \label{fig:median_mad}
\end{figure}
Overall, File 05G demonstrates that:
\begin{itemize}
    \item lightweight robust aggregation can significantly stabilise both FL and SFL under poisoning attacks,
    \item SFL naturally limits the attack surface,
    \item and Median/MAD defence mechanisms are particularly suitable for small-client extreme non-IID environments.
\end{itemize}